\documentclass[../../../include/open-logic-section]{subfiles}

\begin{document}

\olfileid{sth}{story}{predicative}	

\olsection{Predikatif lan Impredikatif}

Himpunan Russell, $R$, ditemtokake liwat $\Setabs{x}{x \notin
x}$. Yen diandharake kanthi luwih jangkep, $R$ iku himpunan
kang anggotane persis kabeh himpunan kang ora dadi anggota
awake dhewe. Dadi nalika nemtokake $R$, kita nguantifikasi
domain kang bakal ngemot $R$ (yen ana).

Iki definisi \emph{impredikatif}. Luwih umum, definisi diarani
impredikatif yen lan mung yen definisi iku nguantifikasi domain
kang ngemot objek kang lagi ditemtokake.
	
Sawise muncul paradoks-paradoks iku, Whitehead, Russell,
Poincar\'{e} lan Weyl nolak definisi impredikatif kaya mangkono
amarga dianggep ``muter kanthi cara kang cacat'':
\begin{quote}
	Analisis tumrap paradoks-paradoks kang kudu diendhani nuduhake
	yen kabeh iku asal saka sawijine lingkaran kang cacat.
	Lingkaran kang dimaksud tuwuh saka pamawas yen sawijining
	koleksi objek bisa ngemot anggota kang mung bisa ditemtokake
	liwat koleksi iku minangka sawijine wutuh[\ldots. \textparagraph]

	Prinsip kang ngidini kita nyingkiri totalitas kang ora sah
	bisa diandharake kaya mangkene: `Apa wae kang nglibatake
	\emph{kabeh} anggota sawijining koleksi ora kena dadi
	salah siji anggota koleksi iku'; utawa, kosok balene:
	`Yen sawijining koleksi, saupama nduweni totalitas,
	bakal nduweni anggota kang mung bisa ditemtokake nganggo
	totalitas iku, koleksi kasebut ora nduweni totalitas.'
	Iki bakal kita arani `prinsip lingkaran cacat', amarga
	prinsip iki ngidini kita nyingkiri lingkaran cacat kang
	tuwuh nalika nganggep anane totalitas kang ora sah.
	\citep[p.~37]{WhiteheadRussell1910}
\end{quote}
Yen kita ngetutake dheweke kanthi nampa \emph{prinsip lingkaran
cacat}, kita bisa nyoba ngganti Skema Komprehensi Naif kang
ngasilake bilai (ing \olref[sth][story][rus]{sec}) nganggo
prinsip kaya mangkene:

\begin{defish}
\emph{Komprehensi Predikatif.} Kanggo saben formula $\phi$ kang
mung nguantifikasi himpunan: himpunan$^\prime$ $\Setabs{x}{\phi(x)}$
ana.
\end{defish}

Anggere himpunan-himpunan$^{\prime}$ dudu himpunan-himpunan
ing tataran kapisan iku, ora ana kontradiksi kang tuwuh.

Sayange, Komprehensi Predikatif ora pati \emph{komprehensif}.
Prinsip iki ngenalake entitas anyar, himpunan-himpunan$^\prime$.
Dadi kita uga kudu nimbang formula kang nguantifikasi
himpunan-himpunan$^\prime$. Yen formula iku tansah ngasilake
himpunan$^\prime$, Paradoks Russell bakal tuwuh maneh nalika
kita nimbang himpunan$^\prime$ saka kabeh
himpunan-himpunan$^\prime$ kang ora dadi anggota awake dhewe.
Mula, miturut gagasan kang padha, kita kudu nganggep formula
kang nguantifikasi himpunan-himpunan$^\prime$ ngasilake
himpunan$^{\prime\prime}$ kang cocog. Banjur kita butuh
himpunan-himpunan$^{\prime\prime\prime}$,
himpunan-himpunan$^{\prime\prime\prime\prime}$, lan sapiturute.
Supaya tandha prima ora saya akeh, luwih gampang yen kita
nganggep iki minangka himpunan$_0$, himpunan$_1$, himpunan$_2$,
himpunan$_3$, himpunan$_4$,\ldots. Iki nggawa kita menyang
teori jinis (prasaja).

Ana sawetara sanggahan kang cetha tumrap teori kaya mangkono
(sanajan durung cetha yen sanggahan iku \emph{nemtokake} banget).
Ringkese: teori kang kawangun angel digunakake; teori iku
ngajokake akeh banget jinis objek kang beda; lan pungkasane
durung cetha yen definisi impredikatif pancen \emph{ala banget}.
	
Kanggo njlentrehake prakara pungkasan iku, gatekna komentar saka
\citeauthor{Ramsey1925} iki:
\begin{quote}
	kita bisa nyebut sawijining priya minangka wong paling dhuwur
	ing sawijining kelompok; kanthi mangkono kita ngenali dheweke
	liwat totalitas kang dheweke dhewe dadi anggotane, tanpa ana
	lingkaran kang cacat. \citep{Ramsey1925}
\end{quote}
Pamawase Ramsey yaiku ``priya paling dhuwur ing kelompok''
\emph{pancen} definisi impredikatif, nanging cetha ora nuwuhake
masalah.

Wong bisa nanggapi yen ing kasus iki kita bisa nemtokake wong
paling dhuwur nganggo cara \emph{predikatif}. Upamane, bisa wae
kita mung nuding priya iku. Mula sanggahan tumrap definisi
impredikatif kudu diwatesi marang entitas kang \emph{mung} bisa
ditemtokake kanthi impredikatif. Nanging malah ing kasus iku,
kita isih butuh katerangan ngapa ``impredikativitas esensial''
iku ala banget.\footnote{Kanggo katrangan luwih, delengen
\citet{Linnebo2010}.}

Pancen, definisi impredikatif dadi masalah gedhe yen kita
ngarepake definisi menehi cara kanggo \emph{nggawe} objek.
Amarga yen definisine impredikatif, kita temenan bakal kajebak
ing lingkaran kang cacat: kanggo nggawe objek kang ditemtokake
kanthi impredikatif, kita \emph{luwih dhisik} kudu nggawe
kabeh objek (kalebu objek kang ditemtokake kanthi impredikatif
iku), amarga objek kasebut ditemtokake nganggo kabeh objek;
dadi kita kudu nggawe objek iku sadurunge objek iku dhewe
digawé. Nanging iki dadi sanggahan serius tumrap himpunan
kang ditemtokake kanthi ``impredikatif sacara esensial'' mung
yen kita nganggep himpunan minangka barang kang kita
\emph{gawe}. Lan kita (mbokmenawa) ora nganggep mangkono.

Mula---becik utawa ala---pendekatan kang banjur lumrah ora
nganut sikap kang kaku tumrap (im)predikativitas. Nanging
pendekatan iki nggunakake apa kang saiki diarani pendekatan
kumulatif-iteratif. Pungkasane, pendekatan iki ngidini kita
merang himpunan dadi ``tataran''---\emph{rada} kaya pendekatan
predikatif merang entitas dadi himpunan$_0$, himpunan$_1$,
himpunan$_2$, \ldots---nanging kita ora ngajokake bedane jinis
antarane himpunan-himpunan iku.

\end{document}
